% !TEX root = ../main.tex
\chapter{Conclusion and Future Work}
\label{ch:conclusion}

Chapter~\ref{ch:discussion} interpreted the findings and judged whether the study met its objectives.

\dissertationSubheading[sec:summary-findings]{Answers to the research questions}

The thesis built a coupled PageRank (C-PR), one walk over the ERC-20 allowance and transfer edges between Arbitrum One contracts, and evaluated it beside its single-layer references, EndorseRank and Adaptive Weighted PageRank (AWP) \parencite{do2023awp}, its ablation S-PR and the raw counts, in a discovery-order test at two freezes and in the holdout windows W0 (April to June~2026) and W1 (July to September~2026). After the registered results were known, the revised analysis was defined: amounts in US dollars on five tokens; scores from the twelve months before each freeze; graphs whose nodes are contracts other than smart accounts, with edges read from the whole chain; each allowance counted up to its owner's balance at the freeze, with balances reconstructed from \texttt{Transfer} logs; constants set on the contract graph; and the discovery-order test. The verdicts below under rules A and B are a re-analysis, and the registered analysis gave the same readings (Appendix~\ref{sec:map-registered}).

RQ1 asked whether EndorseRank and AWP order the same contracts differently, whether each agrees most with current checks of its own edge type, and what each walk costs. They order the cohort differently, and H1 is supported, largely by construction: most cohort contracts hold no capped allowance from another contract and score zero under EndorseRank, so both walks order only about a third of the pairs of contracts (Section~\ref{sec:rank-divergence}). Each walk agrees most with the proxies of its own edges (Table~\ref{tab:alignment-hybrid}), which supports H2 only narrowly, because the proxies come from the same events as the graphs. On the whole contract graph the cost follows the size of the layers a walk covers and the iterations it needs, so EndorseRank, on the far smaller allowance layer, is the cheapest score and C-PR takes slightly longer than AWP (Table~\ref{tab:benchmark-runtime}); H3 is supported, largely by construction.

RQ2 asked whether a crawl ordered by C-PR on the contracts read so far reaches the contracts with the most user-account demand sooner than a breadth-first crawl and than a crawl ordered by the number of contract backlinks. The answer is no. The decision rule of the test's protocol, read at $T_{\mathrm{obs}}$ in the primary reveal mode, gives verdict V2 on both importances (Table~\ref{tab:discovery-decision}): ``Ordered by C-PR on the contracts read so far, a crawl from a random contract of the giant strongly connected component reached the contracts of the universe with the most user-account inflow and those with the most user-authorised value sooner than a breadth-first crawl, but not sooner than a crawl ordered by the number of contract backlinks. A crawl ordered by the number of contract backlinks reached them sooner.'' On user-authorised value its lead over breadth-first order is small and does not hold in every sensitivity (Table~\ref{tab:discovery-sensitivity}). Outside the decision (Table~\ref{tab:discovery-contrasts}), AWP$'$ was the one single-layer walk sooner than the count it smooths on both importances, and the allowance layer added a small, resolved amount to C-PR$'$ over its transfer layer. In the other reveal mode, which shows a contract's out-edges only and carries no verdict, C-PR$'$ was sooner than both comparators at both freezes. Every crawl stayed far from the ideal order (Table~\ref{tab:discovery-auc}). The test decides no rule of the plan or the registration.

RQ3 asked whether scores frozen at a freeze predict later new approval pairs and new transfer senders, and whether any walk, the coupled one included, beats the raw count of the edges it smooths. The first part holds in both windows, and the edge type decides which label a score anticipates (Tables~\ref{tab:holdout-spenders} and~\ref{tab:fresh-holdout}). On the second part the counts mostly win. EndorseRank ranks later approval pairs below the in-approve degree in both windows, with either restart rule (Section~\ref{sec:holdout-results}; Tables~\ref{tab:holdout-diff} and~\ref{tab:fresh-posthoc}). AWP ranks later senders slightly above the transfer in-degree on the whole cohorts, but below it among the W0 spenders that another contract had paid before the freeze (Table~\ref{tab:holdout-receiving}). C-PR ranks later approval pairs above the union in-degree, the count of the edges it walks, and later senders below it; it lies below the in-approve degree on approval pairs and below the transfer in-degree on senders (Table~\ref{tab:cpr-contrasts}). The labels count mostly user accounts, which no score reads, and more than a quarter of each spender cohort lies outside the union graph and scores zero (Tables~\ref{tab:label-split} and~\ref{tab:cohort-coverage}).

RQ4 asked whether the coupled walk ranks contracts better than either of its layers walked alone, in window and out of window. The answer is no. Rule~A, applied to the scores of the contract graph, is not met in either part: in W0 C-PR ranks new approval pairs below its allowance layer and new senders below its transfer layer (Table~\ref{tab:holdout-diff}), and in the same window it lies outside the interval of the leading layer on the transfer and allowance families and does not exceed both layers on Sybil stability (Table~\ref{tab:tau-diff-hybrid}). Among the spenders inside the union graph the contrast with the allowance layer does not exclude zero, a comparison that decides no rule (Table~\ref{tab:union-members}).

RQ5 asked whether adding the allowance layer improves a transfer walk on new data, gaining on later approvals without losing on later inflow. The answer is again no. The thesis applies two of rule~B's registered criteria, F1 and F2; under the registration the rule holds only if each of its criteria holds. F1 holds: in W1 C-PR ranks new approval pairs above its transfer layer walked alone. F2 fails: its loss on new senders exceeds the margin. Rule B is not met, and the two registered sensitivity analyses give the same reading (Table~\ref{tab:fresh-contrasts}). Beside the rules, C-PR ranks later approval pairs below EndorseRank and above AWP and its ablation S-PR, and later senders above EndorseRank and below AWP (Table~\ref{tab:cpr-contrasts}).

Claims~C1--C6 (Section~\ref{sec:empirical-claims}) state these findings formally.

The main research question asked whether one PageRank walk over allowance and transfer edges between Arbitrum One contracts finds the contracts that users rely on sooner, and predicts their later use better, than single-layer walks and raw counts. On these data it does not. Its crawl is sooner than a breadth-first crawl but later than the crawls ordered by the count of contract backlinks and by AWP$'$, it meets neither of the rules that set it against its own layers, and out of window it beats neither single-layer count on the label of that count's edge type. What the allowance layer adds is small and resolved: a little in discovery order, and a gain on later approval pairs that is paid for on later inflow.

\dissertationSubheading[sec:contribution-knowledge]{Original contribution to knowledge}

Read against the literature of Chapter~\ref{ch:literature}, the five contributions of Section~\ref{sec:contributions} establish the following.

In the work reviewed, graph reputation on public ledgers ranked addresses by the transfers among all of them \parencite{do2023awp}, by how risk spreads through transaction networks \parencite{lin2025riskprop}, or by explicit ratings \parencite{bellini2020trust}, and the reputation tooling deployed in DeFi scored the holder of an address from credentials \parencite{humanpassportnddocs}. None of it ranked contracts on the relations among contracts alone, used the ERC-20 allowance \parencite{vogelsteller2015erc20} as the edge of a ranking or coupled it with payments in one walk; approvals had been used to cluster the addresses of one entity \parencite{victor2020clustering} and studied as a security risk \parencite{wang2022unlimited}. None asked of a ledger ranking what the crawl-ordering experiment asked of web rankings \parencite{cho1998crawling}, how soon it leads to what is important, and none set on-chain scores beside the raw counts they smooth, with confidence intervals, paired differences and temporal holdouts. This thesis supplies these.

The new knowledge has five parts. First, a ranking of contracts can be computed from the relations among contracts alone, and what it measures depends on the edges it uses, which can be checked on one cohort with quantified uncertainty. Second, on such a graph the walk that couples allowance and transfer edges leads a crawl to the contracts that users rely on sooner than breadth-first order but not sooner than the plain count of contract backlinks; the result depends on what a visit reveals, and where a visit shows only out-edges the walk leads both, without a verdict. Third, the coupled walk meets neither rule of the plan and the registration: it ranks each label below the layer of that label's edge type walked alone, and over its transfer layer it gains on later approvals but loses more than the margin on later inflow. Fourth, the counts are hard to beat: EndorseRank ranks later approvals below the in-approve degree with either restart rule, C-PR ranks no label above the degree of that label's edge type, and of the single-layer walks only AWP leads its count, in discovery order on both importances and, slightly and on the whole cohorts only, on later inflow. Fifth, an allowance between contracts is rarely backed: counted up to the owner's balance, most such allowances form no edge, and on the thin layer that remains a walk with uniform restarts is at depth one, up to a constant, a linear function of the transition mass into a spender (Section~\ref{sub:theory-closed-form}). The open pipeline, whose published summaries let others recompute every table, is part of the contribution, since automated scores that outsiders cannot inspect are a known concern \parencite{packin2024credit}.

The scope is narrow. The results apply to the contract graph of the five tokens on Arbitrum One in these months, with addresses typed from code read at one block, to the cohorts drawn from it and to the coupling the plan fixed (a fixed $\lambda$, dollar-weighted layers, uniform restarts). The allowance layer between contracts is thin, and at $T_{\mathrm{obs}}$ one owner contract holds most of its pairs, so that contract's allowances weigh on the allowance-side results of that freeze (Table~\ref{tab:graph-layers}). The rule for smart accounts misses some kinds, which stay in the graphs as contracts. No score reads a contract's code or says whether the contract is safe, and none resists Sybils: under uniform restarts a closed farm of contracts collects several times its population share (Section~\ref{sec:sybil-cost-model}).

\dissertationSubheading[sec:future-work]{Future work}

Each suggestion below answers a limitation of Chapter~\ref{ch:discussion}. Any new score, weighting or test it leads to should be fixed before it is scored, replicated on new data and checked against the counts it smooths.

\begin{itemize}
\item Credit decisions about contracts. A ranking that found the contracts users rely on could later serve as one input to credit decisions about contracts. The thesis does not test this use, and its results do not yet support it: the walk did not find those contracts sooner than a count, and such a use would also need an outcome of contracts that bears on repayment, which this study did not have.
\item Importance labels from outside the chain. The discovery-order test measures importance on the chain, as the capped dollars that user accounts pay to a contract or let it spend, and that measure depends on the rule that separates user accounts from contracts. Labels of which contracts matter that are collected outside the chain, by people who examine contracts, would test the orderings against a judgment that shares no data with them.
\item Detection of smart accounts. The rule recognizes Safes and ERC-4337 accounts. The kinds that it misses, and the contracts that services create for a single user, order or payment, stay in the graphs as contracts (Appendix~\ref{sec:data-contract-graph}); a better rule would show how far they shape the transfer layer and the importance measures.
\item A later freeze. The discovery-order test and the comparisons with the counts were defined after the registered results were known, and the two freezes share nine months of their score windows. Repeating the test on a later freeze, with its protocol and its primary reveal mode fixed before the data of that freeze exist, would give a registered answer to RQ2, and other rollups would show whether the answer extends beyond Arbitrum One.
\item A coupled walk with activity restarts. C-PR differs from AWP in its coupling, its weights and its restarts at once, and AWP$'$, which restarts by sending activity, was the soonest walk of the deciding run. A coupled walk whose layers carry the edges and weights of AWP and EndorseRank and whose restarts mix sending and approving activity (C-AWP) would set the coupling alone against the published walk.
\item Adversarial evaluation. The Sybil cost model is analytical, and no attack was carried out on chain. Restricting teleportation and dangling redistribution to a vetted seed set, as TrustRank does \parencite{gyongyi2004trustrank}, would remove the uniform inflow that closed farms exploit, and red-teaming with farms of funded contracts would measure the realized cost of an attack.
\end{itemize}

\dissertationSubheading[sec:closing-reflection]{Closing reflection}

A rollup's contracts are many, and the relations among them are public. A transfer between two contracts records a payment, and an allowance, once it is counted only up to what its owner holds, records a spending right that the owner can back. PageRank will rank either graph, and one walk can follow both. On the data studied here that walk was not the best guide to the contracts that users rely on. A crawl ordered by the plain number of contracts that pay or approve a contract reached them sooner, and out of window the count of each edge type ranked the later arrivals of that type above the coupled walk. For a wallet application, an explorer or an auditor that must choose from public logs which contracts to look at first, the counts on the contract graph are the place to start: they cost least, and here they were hard to beat. A walk earns its cost only where it has been shown to beat its count, as the transfer walk AWP did in discovery order and, narrowly and on the whole cohorts only, on later inflow. The allowance between contracts proved a thin layer whose resolved contributions were small. Whether a ranking built from these relations finds what people who examine contracts would choose is a question for labels that such people supply.
